\section{Introduction}
\label{sec:intro}

\IEEEPARstart{E}{mergency} departments (EDs) differ in the patients they serve and in how they are organized. A tertiary centre receives a larger share of high-acuity and admitted patients than a community hospital, and some EDs see mostly low-acuity patients whose needs resemble those of an urgent care clinic. Across the 16 highest-volume EDs in one Canadian province, the share of patients triaged to level~2 of the Canadian Triage and Acuity Scale (CTAS) ranged from about 3\% to 33\%, and the share at level~4 from about 17\% to 49\% \cite{Holroyd2016_EDCaseMix}. EDs also differ in load, in how long admitted patients wait for an inpatient bed, and in physical layout: some pool all treatment spaces, others run a separate fast track, and others use split (vertical) flow in which ambulatory patients are seen in chairs.

Two operating decisions shape patient flow in every ED. The first is the order in which waiting patients are served: strictly by triage level, by an accumulating priority that grows with waiting time, or by a rule that also uses predicted service time or disposition. The second is streaming: which patients are sent to which treatment area. These two decisions together form what we call a \emph{streaming policy}. ED managers often adopt policies that worked elsewhere, such as a CTAS-based fast track, or that are recommended in reviews. The evidence behind these recommendations comes mostly from single-site studies. Systematic reviews find moderately strong evidence that fast tracks shorten waiting times and length of stay (LOS) but only limited evidence for other forms of streaming \cite{Oredsson2011_TriageFlowReview}, and operations research studies show that the value of streaming depends on the setting \cite{Saghafian2012_Streaming_as_a_Mechanism,laker2014flex_track,Hu2009}.

This is a question of fit, in the sense of contingency research in operations management: the value of a practice depends on the context in which it is used \cite{Sousa2008_ContingencyOM}. For ED streaming the question has two parts: which policy fits which type of ED, and how much an ED loses when it uses a policy that does not fit. The second part matters to health system managers who standardize practice across many sites. The cost of a mismatch should be known before a region-wide rollout.

A second, more technical question arises because streaming policies are increasingly implemented in decision support software as scoring rules. A scoring rule ranks patients by a weighted combination of coded attributes, for example by their similarity to a target profile \cite{OFFODILE1990,GT-in-ED_HCMS2016}. Managers and clinicians usually describe the intended policy in words (``acuity first, then quick patients''), and the weights are chosen to express that intent. Whether a given set of weights actually applies the intended rule is rarely checked. If it does not, an evaluation of the ``policy'' is really an evaluation of something else.

A third question follows from the second. If a scoring rule can express many policies, its weights and targets can be chosen for a particular ED rather than by judgement. How much does such tailoring gain over the best of the classical rules, and can a rule tailored for one ED be used in another?

This paper addresses five research questions:
\begin{itemize}
\item[RQ1:] Does the best streaming policy differ across ED types defined by acuity mix, load and configuration?
\item[RQ2:] How much does an ED lose, relative to the best policy for its type, when one policy is used in every ED type?
\item[RQ3:] How accurate must predictions of patient attributes be before prediction-based policies outperform policies that use triage information only?
\item[RQ4:] Under which conditions does a similarity-based scoring rule apply a classical rule exactly, and how close is it otherwise, at the level of individual decisions and of outcomes?
\item[RQ5:] How much can a similarity rule whose parameters are optimized for an ED type improve on the best classical policy, and how well does an optimized rule transfer to other ED types?
\end{itemize}

We answer these questions with a discrete-event simulation of \NEdTypes{} synthetic ED types (three acuity mixes, three load levels and three configurations), \NQueueRules{} queue rules and three routing rules. Patient pathways are drawn from a published process-mining study of ED patient journeys \cite{Rismanchian2023_PatientJourney}; case-mix and admission parameters come from published Canadian aggregates \cite{Holroyd2016_EDCaseMix}. Predictions are not produced by a machine learning model but simulated: each patient has true and predicted attributes, and the predicted attributes are derived from the true ones through confusion matrices whose error is scaled by a single parameter. This lets us vary prediction quality continuously, including perfect information. The parameters of the similarity rules are optimized by a cross-entropy search with separate training, validation and test replications. The conclusions are tested on misspecified inputs and on \NOosTypes{} further ED types that were not used to derive them.

The paper makes five contributions. First, it compares streaming policies across a designed set of ED types under common random numbers and reports, for each ED type, the best policy and the set of policies that cannot be distinguished from it statistically. Second, it measures the regret of one-size-fits-all policies, including the common practice of CTAS ordering with a CTAS-based fast track, and turns the results into a simple decision aid. Third, it reports how much prediction accuracy is needed for prediction-based policies to pay off in each ED type. Fourth, it states conditions under which a similarity-based rule reproduces a classical rule exactly, and measures fidelity when the conditions fail. The fidelity analysis shows that two rules can give almost the same average outcomes while choosing different patients at a large share of decisions, so outcome comparisons alone do not show that a rule has been implemented as intended. Fifth, it optimizes the parameters of the similarity rule for each ED type and objective, measures the gain over the best classical policy and the loss when an optimized rule is moved to another ED type, and shows that the policy class contains all the classical rules studied, so the optimized rule can be read against them. All results are produced by one open-source code base in which every number in the paper is generated from the simulation output.

The rest of the paper is organized as follows. Section~\ref{sec:lit} reviews related work. Section~\ref{sec:framework} describes the simulation framework, the policies and the fidelity propositions. Section~\ref{sec:optimization} describes the optimization of similarity rules. Section~\ref{sec:design} sets out the experimental design and performance measures. Section~\ref{sec:results} reports the results, Section~\ref{sec:implications} discusses managerial implications, Section~\ref{sec:limitations} discusses limitations, and Section~\ref{sec:conclusion} concludes. The online appendix gives the attribute coding, the proofs, the full model inputs, the algorithms and detailed results, so that the study can be reproduced from the paper and the appendix alone.
